\section{実験結果と考察}
\label{sec:experiments}

\subsection{データ生成方法と実験プロトコル}
\label{sec:exp_method}

本研究では，バスケット構造を持つ合成トランザクションデータを生成し，
提案手法と従来 Apriori-window（flatten）を同一条件で評価した．
トランザクション $t$ のバスケット数 $N_t$ と，バスケット $b$ のアイテム数 $K_{t,b}$ は
次で与える．
\[
N_t = \max\{1, \mathrm{Poisson}(\lambda_{\text{baskets}})\},
\]
\[
K_{t,b} = \max\{1, \mathrm{Poisson}(\lambda_{\text{basket-size}})\}.
\]
アイテムはカテゴリ集合
\[
C_g = \{\,i \mid i \bmod G = g\,\}
\]
で分割し，各バスケットでは主カテゴリを一様選択した後，
確率 $p_{\text{same}}$ で同カテゴリ，確率 $1-p_{\text{same}}$ で全体集合から
Zipf 重み $w_i \propto (i+1)^{-\alpha_{\text{zipf}}}$ に従って抽出する．

評価指標は，E0 では $(P,\text{window-start})$ 単位の FP/FN，
E1/E2 では偽共起率（GTSPR），
E3 では実行時間と Trad./BAAW 比（従来法時間 / Basket-aware Apriori-Window 時間），
E4 では区間 purity，GT 区間との mean Jaccard，
従来法の spurious interval rate，および non-spurious expansion である．
E2 では補助指標として observed spurious（従来法のみが検出した多項目パターン数）も報告する．
E1/E2/E4 は seed=5，E3 は seed=3 を用い，表では seed 平均を報告する．
詳細な固定パラメータと sweep 設定は補遺（Sec.~\ref{sec:supp_exp_conditions}）に示す．

\subsection{E0: 実装厳密性の検証（Oracle 一致）}
\label{sec:exp_e0}

E0 の目的は，Basket-aware Apriori-Window 実装が定義どおりに動作していることを，
総当たり Oracle との一致で確認することである．
具体的には，小規模合成データ上で
\begin{itemize}
    \item 全アイテム集合候補（$|P| \le L$）と全ウィンドウ開始位置を総当たりで評価した Oracle 出力
    \item Basket-aware Apriori-Window の出力
\end{itemize}
を比較し，$(P,\text{window-start})$ 単位で FP/FN を計測した．

実験は 34 ケース（パラメータ組合せ）で実施し，結果は全ケースで
FP=0, FN=0 であった（表~\ref{tab:e0_summary}）．
この結果は，後続比較（E1）の前提として，実装バグによる過検出・見落としが
観測されなかったことを示す．

\begin{table}[t]
\caption{E0 結果（Oracle 一致）}
\label{tab:e0_summary}
\centering
\begin{tabular}{lc}
\toprule
項目 & 値 \\
\midrule
評価ケース数 & 34 \\
PASS ケース数 & 34 \\
FAIL ケース数 & 0 \\
合計 FP & 0 \\
合計 FN & 0 \\
\bottomrule
\end{tabular}
\end{table}

\subsection{E1: $\lambda_{\text{baskets}}$ sweep による偽共起拡大}
\label{sec:exp_e1}

E1 では，$\lambda_{\text{baskets}} \in \{1.0,1.5,2.0,3.0,5.0\}$ を sweep し，
偽共起の拡大傾向を評価した（seed=5）．
比較対象は Basket-aware Apriori-Window と従来 Apriori-window（flatten）である．

まず予備確認として，従来法の GTSPR が $\lambda_{\text{baskets}}$ の増加とともに
単調増加し，Basket-aware Apriori-Window が全条件で 0 を維持することを確認した
（表~\ref{tab:prelim_gtspr}）．
従来法は 0.454 から 0.971 まで増加し，バスケット密度の上昇に伴って
偽共起が支配的になることがわかる．

\begin{table}[t]
\caption{E1 予備結果: GTSPR（seed=5 平均）}
\label{tab:prelim_gtspr}
\centering
\begin{tabular}{ccc}
\toprule
$\lambda_{\text{baskets}}$ & Apriori-window (flatten) GTSPR & Basket-aware Apriori-Window GTSPR \\
\midrule
1.0 & 0.454 & 0.000 \\
1.5 & 0.676 & 0.000 \\
2.0 & 0.794 & 0.000 \\
3.0 & 0.906 & 0.000 \\
5.0 & 0.971 & 0.000 \\
\bottomrule
\end{tabular}
\end{table}

\subsection{E2: 構造感度（$G$ / $p_{\text{same}}$ sweep）}
\label{sec:exp_e2}

E2 では，データ構造に対する偽共起感度を確認するため，
2 種類の sweep を実施した（seed=5）．
すなわち，
\begin{itemize}
    \item $G \in \{3,5,10,20,50\}$（$p_{\text{same}}=0.8$ 固定）
    \item $p_{\text{same}} \in \{0.5,0.6,0.7,0.8,0.9\}$（$G=10$ 固定）
\end{itemize}
である．

表~\ref{tab:e2_g_sweep} より，従来法 GTSPR は
$G=3$ の 0.827 から $G=50$ の 0.638 まで単調減少した．
Basket-aware Apriori-Window は全条件で GTSPR=0 を維持した．
また observed spurious も 24.8 から 3.0 に減少し，
カテゴリ粒度が粗い条件で偽共起リスクが高いことが確認できる．

表~\ref{tab:e2_psame_sweep} より，従来法 GTSPR は
$p_{\text{same}}=0.5$ の 0.827 から $0.9$ の 0.741 まで単調減少した．
Basket-aware Apriori-Window は同様に全条件で 0 であった．
observed spurious も 19.8 から 5.4 に減少した．

\begin{table}[t]
\caption{E2-G 結果（$p_{\text{same}}=0.8$, seed=5 平均）}
\label{tab:e2_g_sweep}
\centering
\begin{tabular}{cccc}
\toprule
$G$ & Apriori-window GTSPR & Basket-aware Apriori-Window GTSPR & Obs. Spur. \\
\midrule
3  & 0.827 & 0.000 & 24.8 \\
5  & 0.822 & 0.000 & 19.0 \\
10 & 0.794 & 0.000 & 10.6 \\
20 & 0.722 & 0.000 & 5.8 \\
50 & 0.638 & 0.000 & 3.0 \\
\bottomrule
\end{tabular}
\end{table}

\begin{table}[t]
\caption{E2-$p_{\text{same}}$ 結果（$G=10$, seed=5 平均）}
\label{tab:e2_psame_sweep}
\centering
\begin{tabular}{cccc}
\toprule
$p_{\text{same}}$ & Apriori-window GTSPR & Basket-aware Apriori-Window GTSPR & Obs. Spur. \\
\midrule
0.5 & 0.827 & 0.000 & 19.8 \\
0.6 & 0.822 & 0.000 & 17.2 \\
0.7 & 0.814 & 0.000 & 15.2 \\
0.8 & 0.794 & 0.000 & 10.6 \\
0.9 & 0.741 & 0.000 & 5.4 \\
\bottomrule
\end{tabular}
\end{table}

\subsection{E3: スケーラビリティ（$N_{\text{transactions}}$ sweep）}
\label{sec:exp_e3}

E3 では，$N_{\text{transactions}} \in \{10^3,10^4,10^5,10^6\}$ を sweep し，
大規模条件での実行時間スケーリングを評価した（seed=3）．
他パラメータは E1 と同一である．

表~\ref{tab:e3_scalability} より，両手法の実行時間は $N$ とともに増加し，
従来法/提案法の時間比（Trad./BAAW）も
0.945，1.076，1.161，1.345 と単調増加した．
最小規模（$10^3$）では差は小さいが，
$10^4$ 以降は一貫して Basket-aware Apriori-Window が高速であり，
$10^6$ では Basket-aware Apriori-Window が 108.5 秒，従来法が 145.9 秒（1.345 倍）であった．
この結果は，規模拡大時に提案法の計算優位が増大することを示す．

\begin{table}[t]
\caption{E3 結果（実行時間, ms; seed=3 平均）}
\label{tab:e3_scalability}
\centering
\begin{tabular}{cccc}
\toprule
$N_{\text{transactions}}$ & Basket-aware Apriori-Window & Traditional & Trad./BAAW \\
\midrule
$10^3$ & 15.2 & 14.3 & 0.945 \\
$10^4$ & 177.7 & 191.2 & 1.076 \\
$10^5$ & 2962.1 & 3435.0 & 1.161 \\
$10^6$ & 108506.3 & 145851.9 & 1.345 \\
\bottomrule
\end{tabular}
\end{table}

\subsection{E4: 区間品質評価（$\lambda_{\text{baskets}}$ sweep）}
\label{sec:exp_e4}

E4 では，E1 と同一の $\lambda_{\text{baskets}} \in \{1.0,1.5,2.0,3.0,5.0\}$（seed=5）で，
検出された密集区間そのものの品質を評価した．
具体的には，以下の 2 点を確認する．
\begin{itemize}
    \item 検出区間内で真の同一バスケット共起が占める割合（interval purity）
    \item 各検出区間と対応 Ground Truth 区間の重なり（mean Jaccard）
\end{itemize}
加えて従来法について，
同一 itemset に対して Basket-aware Apriori-Window 区間と重ならない区間の割合
（spurious interval rate）と，
非 spurious itemset 上での区間長増分（non-spurious expansion）を計測した．

表~\ref{tab:e4_interval_quality} より，Basket-aware Apriori-Window は全 $\lambda$ で
従来法より高 purity・高 Jaccard を示した．
とくに mean Jaccard は全条件で 1.000 であり，
Ground Truth 区間境界と一致することを確認した．
一方，従来法の purity は $\lambda=1.0$ の 0.075 から
$\lambda=5.0$ の 0.018 まで単調減少し，
高密度条件ほど区間汚染が進むことが分かる．

表~\ref{tab:e4_interval_artifacts} より，従来法の spurious interval rate は
0.735--0.945 と高水準で推移し，
$\lambda=5.0$ で最大となった．
さらに non-spurious expansion は全条件で正であり，
偽共起のみならず，真の共起区間に対しても境界が過剰に拡張される傾向を確認した．
この結果は，E1 の集合レベル誤検出に加えて，
時間区間レベルでも提案法の有効性があることを示す．

\begin{table}[t]
\caption{E4 結果（区間品質; seed=5 平均）}
\label{tab:e4_interval_quality}
\centering
\scriptsize
\begin{tabular}{ccccc}
\toprule
$\lambda$ & BAAW purity & Trad. purity & BAAW Jaccard & Trad. Jaccard \\
\midrule
1.0 & 0.096 & 0.075 & 1.000 & 0.098 \\
1.5 & 0.096 & 0.061 & 1.000 & 0.045 \\
2.0 & 0.095 & 0.054 & 1.000 & 0.097 \\
3.0 & 0.096 & 0.043 & 1.000 & 0.141 \\
5.0 & 0.102 & 0.018 & 1.000 & 0.035 \\
\bottomrule
\end{tabular}
\end{table}

\begin{table}[t]
\caption{E4 結果（従来法の区間汚染・膨張; seed=5 平均）}
\label{tab:e4_interval_artifacts}
\centering
\scriptsize
\begin{tabular}{cccc}
\toprule
$\lambda$ & Trad. spur. rate & Trad. expansion & Trad. exp.$>0$ ratio \\
\midrule
1.0 & 0.735 & 2064.7 & 1.000 \\
1.5 & 0.834 & 6129.0 & 1.000 \\
2.0 & 0.856 & 2543.5 & 0.778 \\
3.0 & 0.803 & 1371.0 & 0.858 \\
5.0 & 0.945 & 1755.7 & 0.734 \\
\bottomrule
\end{tabular}
\end{table}

\subsection{E6: 実データ検証（dunnhumby / onlineretail / chicago）}
\label{sec:exp_e6}

E6 では，合成データで観測した傾向の外的妥当性を確認するため，
3 つの実データで比較を行った．
入力は時刻付き生データから再構築したバスケット構造付き系列であり，
集約幅はバスケット過密を抑えるように選定した
（dunnhumby: 1h，onlineretail: 3h，chicago: 30m）．
最終的なトランザクション数はそれぞれ 15,240 / 2,338 / 297,777 である．

主比較（E6-1）は，各データセットで固定した baseline 設定
（dunnhumby: $W=336,\theta=60,L=3$，
onlineretail: $W=72,\theta=40,L=3$，
chicago: $W=48,\theta=20,L=2$）で
Basket-aware Apriori-Window（BAAW）と従来法（flatten）を比較した．
結果を表~\ref{tab:e6_real_main} に示す．
3 データセットすべてで Observed Spurious が正であり，
従来法検出多項目パターン中に占める割合は
0.742--0.992（平均 0.906）であった．
特に onlineretail では 132,761 パターン中 131,702 が
BAAW では検出されず，flatten による構造混入が支配的であることが確認できる．

\begin{table}[t]
\caption{E6-1 結果（実データ baseline 比較）}
\label{tab:e6_real_main}
\centering
\scriptsize
\begin{tabular}{lrrrrr}
\toprule
Dataset & BAAW multi & Trad. multi & Obs. Spur. & Spur. rate & Trad./BAAW time \\
\midrule
dunnhumby    & 37   & 2,418   & 2,383   & 0.986 & 1.012 \\
onlineretail & 1,152 & 132,761 & 131,702 & 0.992 & 4.521 \\
chicago      & 670  & 2,593   & 1,923   & 0.742 & 0.750 \\
\bottomrule
\end{tabular}
\end{table}

E6-2（感度分析）では，dataset ごとに小規模 sweep を行った．
dunnhumby では $\theta$ を 60 から 50 に下げると
Trad. 多項目パターン数が 2,418 から 5,463 に増加し，
BAAW（37 $\rightarrow$ 60）より増分が大きかった．
chicago でも $W=24 \rightarrow 96$（$\theta=20$）で
BAAW 492 $\rightarrow$ 934 に対して Trad. 1,801 $\rightarrow$ 3,517 と，
常に Trad. が高い件数を示した．
onlineretail では $\theta=40$ で差分が最大（1,152 vs 132,761）となり，
$\theta=60$ では双方ほぼ 0 に近づくため，偽共起の顕在化領域は
閾値設定に依存することが分かる．

E6-3（規模変化）では，30d/180d/365d/all の horizon で実行時間を比較した．
全 dataset・全手法で runtime は horizon 拡大に対して単調増加し，
特に onlineretail では Trad./BAAW 比が
0.912（30d）$\rightarrow$ 2.839（180d）$\rightarrow$ 4.600（365d）へ増加した．
一方，dunnhumby/chicago では本設定で Trad. が高速な領域も存在したが，
表~\ref{tab:e6_real_main} の通り品質面では Trad. の偽共起混入率が高く，
速度のみでは手法選択できないことを示す．

\subsection{考察}
\label{sec:discussion}

\subsubsection{統計的妥当性}
E1/E2/E4 は seed 5 本，E3 は seed 3 本の平均を報告した．
BAAW の GTSPR は全 seed・全条件で厳密に 0 であり，
これは確率的結果ではなく定義から導かれる決定論的性質である（命題~\ref{prop:spr_zero}）．
一方，従来法の GTSPR は seed 間分散が小さく（全条件で変動係数 $< 0.05$），
報告した単調傾向は統計的に頑健である．

\subsubsection{BAAW の purity 値について}
表~\ref{tab:e4_interval_quality} の BAAW purity が約 0.096 と低い値を示すが，
これは purity の定義（区間内全トランザクション中，対象パターンが共起するトランザクションの割合）に
起因する．密集区間の最小サポート $\theta$ を下回らない保証はあるが，
区間内の全トランザクションでパターンが出現するわけではないため，
purity は $\theta / N_{\text{window}}$ 程度の値となるのが正常な挙動である．

\subsubsection{実データにおける偽共起の実務的影響}
E6 の結果は，集約粒度の選択が分析結果に決定的な影響を与えることを示唆する．
特に onlineretail では 99.2\% の多項目パターンが偽共起であり，
これらに基づく推薦や需要予測は誤った共起関係に立脚することになる．
BAAW はこの問題をアルゴリズム側で構造的に解決し，
データ前処理への依存を排除する．
